\begin{tabular}{cr}\toprule
$j$ & $\displaystyle\int_j^{j+1}R(x)x^{-3}\,dx$\\\midrule
3 & $\frac{26807}{161280}$\\
4 & $\frac{37383}{704000}$\\
5 & $\frac{37523}{8236800}$\\
6 & $-\frac{120923}{6552000}$\\
7 & $-\frac{10025233}{356428800}$\\
8 & $-\frac{110029309}{3348864000}$\\
9 & $-\frac{2278419487}{64465632000}$\\
10 & $-\frac{282415081}{7724640000}$\\
11 & $-\frac{3236921227}{87524236800}$\\
12 & $-\frac{3350220001}{90899827200}$\\
13 & $\frac{7424224373}{2217983040000}$\\
14 & $\frac{16775764609}{955086612480}$\\
15 & $\frac{369043847}{30810528000}$\\
16 & $\frac{651380108633}{86461373856000}$\\
17 & $\frac{472851276229}{119820121344000}$\\
18 & $\frac{4930060867}{4724197793280}$\\
19 & $-\frac{15199801}{11563552000}$\\
\bottomrule\end{tabular}
\end{table}

\subsection{The outer integral}
Let
\[
 T(y)=R_0(y)-d(y)-2\lambda\floor{1/y}+\sum_{j=1}^5(2\lambda-jy)_+.
\]
On each open interval of Table~\ref{tab:outerintegrals}, $T(y)=b+cy$.
Resolving every minimum, positive part, and reciprocal floor by exact
endpoint comparisons proves those affine expressions. Their integrals
are $b(r-l)+c(r^2-l^2)/2$, and their sum is
$I_{\mathrm{out}}=127751/96000$. As before, no assertion about an affine
branch is inferred from numerical sampling.
\gid{Zeta5.Certificate.outer_integral}
\begin{table}[htbp]\centering\small
\caption{Outer integrand $T(y)=b+cy$ (source Table 4).}\label{tab:outerintegrals}
\renewcommand{\arraystretch}{1.65}
\begin{tabular}{rrrr}\toprule
$l$ & $r$ & $b$ & $c$\\\midrule
$\frac{1}{3}$ & $\frac{43}{120}$ & $\frac{279}{40}$ & $-9$\\
$\frac{43}{120}$ & $\frac{37}{100}$ & $\frac{451}{40}$ & $-21$\\
$\frac{37}{100}$ & $\frac{13}{30}$ & $\frac{377}{40}$ & $-16$\\
$\frac{13}{30}$ & $\frac{37}{80}$ & $\frac{429}{40}$ & $-19$\\
$\frac{37}{80}$ & $\frac{1}{2}$ & $\frac{17}{4}$ & $-5$\\
$\frac{1}{2}$ & $\frac{43}{80}$ & $\frac{51}{10}$ & $-3$\\
$\frac{43}{80}$ & $\frac{37}{60}$ & $\frac{231}{20}$ & $-15$\\
$\frac{37}{60}$ & $\frac{13}{20}$ & $\frac{97}{10}$ & $-12$\\
$\frac{13}{20}$ & $\frac{37}{40}$ & $\frac{42}{5}$ & $-10$\\
$\frac{37}{40}$ & $1$ & $1$ & $-2$\\
$1$ & $\frac{37}{20}$ & $\frac{37}{20}$ & $-1$\\
\bottomrule\end{tabular}
\end{table}

\subsection{The final rational margins}
Direct substitution into \eqref{eq:AM} gives
\begin{align*}
 A_{200}&=\frac{127125602969131786927559}{94195881588024216000000},\\
 A_{100000}&=\frac{7756864096839411316755964319057}
                  {5887242599251513500000000000000}.
\end{align*}
With $U=-2733991/2000000$, the exact differences are
\begin{align}
 -1600(A_{200}+U)-\frac{139}{5}
 &=\frac{3089837638249482469}{58872425992515135000}>0,
 \label{eq:margin200}\\
 -1600(A_{100000}+U)-\frac{7907}{100}
 &=\frac{29873543950273155160680943}
          {3679526624532195937500000000}>0.
 \label{eq:marginlarge}
\end{align}
These imply all strict inequalities used in Section~\ref{sec:finish}.
In particular $A_{100000}+U<-79/1600$.
\gid{Zeta5.Certificate.final_rational_margins}

\section{Optional rational approximation and coefficient height}
\label{app:approx}
This appendix retains source Appendix C. Its matrix norm bound is not an
assumption of Theorem~\ref{thm:main}.

\subsection{Relative norm of the coefficient matrix}
Put
\[
 G=G_K(\xi),\qquad M_K=[X]G_K(X),\qquad
 B_K=G^{-1/2}M_KG^{-1/2}.
\]
The square root and inverse are the real symmetric positive-definite
ones supplied by Proposition~\ref{prop:moments}. Norm means Euclidean
operator norm.

\begin{proposition}[Relative norm; source Proposition C.1]
\label{prop:relativenorm}
For all sufficiently large $K\in40\Z_{>0}$,
\begin{equation}\label{eq:relativenorm}
 \|B_K\|\le e^{\varrho K},\qquad\varrho=\frac{6933}{500}.
\end{equation}
\gid{Zeta5.Approx.relative_matrix_norm}
\end{proposition}
\begin{proof}
Let $q$ be a real polynomial of degree less than $h$ and put
$p(t)=q(K^2t)$ and $A_K=K^{12N-2K+6}$. For any $0<a<b$, keep only the
$\ell=1$ term of the positive weight, restrict $y=K\sqrt t$ to
$a\le t\le b$, and bound every factor separately. This gives
\begin{equation}\label{eq:Glower}
 \boldsymbol q^{\mathsf T}G\boldsymbol q
 \ge A_Kc_a e^{-F(a,b)K}\int_a^b p(t)^2\,dt,
\end{equation}
where
\[
 c_a=\frac{(2\pi)^4a^2}{24},\qquad
 F(a,b)=2\pi\sqrt b+\log(1+b)-6\alpha\log a.
\]
The scaled numerator is at least $a^{6N}$, the denominator is at most
$(1+b)^K$, and the remaining scaled density is at least
$c_a e^{-2\pi K\sqrt b}$. This proves the bound with its exact power
of $K$.

Write $u_j=j/K$. The residue expression in
Proposition~\ref{prop:degree} gives
\begin{equation}\label{eq:Mresidue}
 \boldsymbol q^{\mathsf T}M_K\boldsymbol q
 =A_K\sum_{j=N+1}^K r_jp(-u_j^2)^2,
 \quad
 r_j=\frac{u_j^4\prod_{\ell=1}^N(u_\ell^2-u_j^2)^5}
               {\prod_{N<k\le K,\ k\ne j}(u_k^2-u_j^2)}.
\end{equation}
The numerator has absolute value at most one. Since $u_k+u_j\ge2\alpha$,
\begin{align*}
 |r_j|&\le
 \frac{(K/(2\alpha))^{h-1}}{(j-N-1)!(K-j)!}
 \le\frac{(K/\alpha)^{h-1}}{(h-1)!}\\
 &\le\exp\{\lambda K(1-\log(\alpha\lambda))\}.
\end{align*}
The middle inequality is the bound $\binom{h-1}{j-N-1}\le2^{h-1}$.
For the last one, use $m!\ge(m/e)^m$ and monotonicity of
$x\log(eK/(\alpha x))$ for $0<x\le h$.

Map $[a,b]$ to $[-1,1]$ and expand $p$ in the orthogonal Legendre basis.
For $-1\le t\le0$ the transformed argument has absolute value at most
\[
 Y=\frac{2+a+b}{b-a}>1,
 \qquad R=Y+\sqrt{Y^2-1}.
\]
The Laplace integral for the Legendre polynomial gives
$|P_j(s)|\le R^j$ on this exterior interval (use parity for negative
$s$). Cauchy--Schwarz and $\sum_{j=0}^{h-1}(2j+1)=h^2$ imply
\begin{equation}\label{eq:legendreexterior}
 \max_{-1\le t\le0}|p(t)|^2
 \le\frac{h^2R^{2h}}{b-a}\int_a^b p(t)^2\,dt.
\end{equation}
Combining \eqref{eq:Glower}--\eqref{eq:legendreexterior} and using the
Rayleigh quotient for the real symmetric matrix $B_K$ gives
\[
 \|B_K\|\le\frac{h^3}{c_a(b-a)}e^{\kappa(a,b)K},\qquad
 \kappa(a,b)=F(a,b)+\lambda(1-\log(\alpha\lambda))+2\lambda\log R.
\]
Choose $a=1/18$, $b=1/3$, so $Y=43/5$ and
\begin{equation}\label{eq:kappaexplicit}
\begin{split}
 \kappa={}&\frac{2\pi}{\sqrt3}+\log\frac43+\frac9{20}\log18
 +\frac{37}{40}\left(1+\log\frac{1600}{111}\right)\\
 &+\frac{37}{20}\log\frac{43+4\sqrt{114}}5
 <\frac{6933}{500}.
\end{split}
\end{equation}
The strict margin is certified in \eqref{eq:kappacert}. Since a fixed
polynomial factor in $h$ is eventually at most
$e^{(\varrho-\kappa)K}$, it is absorbed by this positive margin. This
proves \eqref{eq:relativenorm}.
\end{proof}

\subsection{The approximation contradiction with explicit quantifiers}
\begin{proof}[Proof of Corollary~\ref{cor:approx}]
Fix $M=100000$ and write $Q_K=Q_{K,M}$. From
\eqref{eq:finalasymptotic} and Appendix~\ref{app:arithmetic}, for all
sufficiently large allowed $K$,
\[
 0<Q_K(\xi)\le e^{-\epsilon K^2},\qquad\epsilon=79/1600.
\]
Let $c=75/4$. For an integer $b\ge2$, choose
\[
 K=40\left\lceil\frac{c\log b}{40}\right\rceil,
 \qquad c\log b\le K<c\log b+40.
\]
As $b\to\infty$, this $K$ satisfies admissibility and all eventual
bounds above. Suppose $a\in\Z$ and $\delta=a/b-\xi$ satisfies
$|\delta|\le b^{-260}$. The exact rational identities
\[
 \varrho c-260=-1/80,\qquad\epsilon-\lambda/c=1/24000
\]
give
$|\delta|\|B_K\|\le e^{40\varrho}b^{-1/80}=o(K^{-1})$.
All eigenvalues of $I+\delta B_K$ are therefore positive for sufficiently
large $b$, uniformly in $a$, and
\[
 0<\det(I+\delta B_K)
 \le(1+|\delta|\|B_K\|)^h
 \le e^{h|\delta|\|B_K\|}\longrightarrow1.
\]
It is eventually at most two. Affine dependence on $X$ gives the exact
identity
$Q_K(a/b)=Q_K(\xi)\det(I+\delta B_K)$. Integrality and degree $h$ now
imply that $b^hQ_K(a/b)$ is a positive integer, whereas
\[
 b^hQ_K(a/b)\le2\exp(\lambda K\log b-\epsilon K^2)
              \le2e^{-K^2/24000}<1
\]
for every sufficiently large $b$. This contradiction proves the
corollary with a single threshold $b_0$ independent of $a$.
\end{proof}

\subsection{Coefficient height and primitive normalization}
If $H(Q)$ denotes the largest absolute value of an integer coefficient,
the real eigenvalues $\beta_1,\ldots,\beta_h$ of $B_K$ give
\[
 Q_K(X)=Q_K(\xi)\prod_{j=1}^h(1+(X-\xi)\beta_j).
\]
Since $\xi<2$, the product of the coefficient $\ell^1$ norms bounds the
height, and Proposition~\ref{prop:relativenorm} gives
\begin{equation}\label{eq:height}
 \log H(Q_K)\le(\lambda\varrho-\epsilon)K^2+o(K^2)
 =\frac{511067}{40000}K^2+o(K^2)<13K^2
\end{equation}
for sufficiently large allowed $K$. For instance, each factor has
coefficient sum at most $1+3e^{\varrho K}$, which proves the required
$O(K)$ remainder after multiplication by $h$.
\gid{Zeta5.Approx.coefficient_height}

For a nonzero $A\in\Q[X]$, let $\cont_{\Q}(A)>0$ be its rational
content: $A/\cont_{\Q}(A)$ is integral and primitive. Set
\[
 P_K=\Delta_K/\cont_{\Q}(\Delta_K).
\]
For every admissible $K,M$ there is an integer $a_{K,M}\ge1$ with
$Q_{K,M}=a_{K,M}P_K$. To see this, a rational multiple of a primitive
integer polynomial is integral only if the scalar is an integer, by
B\'ezout applied to its coefficients. Here the scalar is positive.
Thus $0<P_K(\xi)\le Q_{K,M}(\xi)$, and the established upper bounds
also hold for the primitive polynomial. The empirical primitive values
reported in the supplied audits are not needed for this deduction.
\gid{Zeta5.primitive_normalization}

\section{Autonomous Lean formalization specification}
\label{app:lean}
\subsection{The endpoint and the trust boundary}
The accompanying declaration ledger is an implementation blueprint, not
a list of already formalized results. At the time of this edition no Lean
proofs in this package have been checked. The finite Python verification
has been executed successfully, but it certifies only the displayed
finite arithmetic and interval enclosures, not the asymptotic or local
functional theorems.

Use a named real definition
\[
 \texttt{zetaFive}:=\sum_{n=0}^{\infty}\frac1{(n+1)^5}
\]
with real coercions and a summability proof. The final unconditional
statement must be \texttt{Irrational zetaFive}, together with a theorem
identifying it with the intended $\zeta(5)$ value. A theorem about an
abstract real parameter satisfying the polynomial estimates is the
useful conditional endpoint, not a substitute for this theorem.

During construction, a structure of named proof inputs is acceptable.
Its fields must be explicit propositions, not project axioms, and the
conditional theorem must retain that structure as a visible argument.
Standard external theorems must likewise be identified by exact type
and provenance. An unconditional endpoint may be exported only after
every input has been instantiated by checked proofs or actual imported
theorems with audited dependencies. No incomplete definition,
\texttt{sorry}, \texttt{admit}, false premise, or added axiom may silently
close an obligation.

The final certificate path must be kernel-checkable. An external Python
program may emit rational witnesses and subdivision data; the theorem
asserting that they imply the desired inequalities must be proved in
Lean. Unchecked numerical or native execution is not by itself a proof
of a certificate. Record any additional computational trust dependency
rather than hiding it. Inspect \texttt{\#print axioms} for every public
endpoint and retain the output in the build audit; see the Lean reference
manual for axiom tracking and validation \cite{LeanValidation}.

\subsection{Minimal mathematical interfaces for external inputs}
The following interfaces state the mathematical content to be found in
libraries or proved; they do not guess current declaration names.

\begin{description}[leftmargin=1.3cm,style=nextline]
\item[Bernoulli arithmetic.]
$B_1=-1/2$; the polynomial difference, reflection, and multiplication
identities; $B_{2k+1}=0$ for $k\ge1$; and, for primes $p$ and $m\ge1$,
$v_p(B_{2m})\ge-1$ with $p$-integrality if $p-1\nmid2m$. The
von Staudt--Clausen denominator statement is a sufficient interface.

\item[Even zeta and Hermite.]
Euler's formula at positive even integers, the identification
$\zeta(5,j)=\xi-\Hfive{j-1}$ for integers $j\ge1$, and, for $a>0$,
\[
 \zeta(5,a)=\frac1{2a^5}+\frac1{4a^4}
 +2\int_0^\infty\frac{\operatorname{Im}(a-iy)^{-5}}
                          {e^{2\pi y}-1}\,dy,
\]
including absolute integrability of the integral. An independently
